\scanpage{12}
\mspage{1}
Espaces stratifiés (étude par voisinages \deleted{tubulaires}
\inserted{coniques})
\textcircled{I}

\textcircled{1} Soit $Y$ un espace topologique. Un \uline{voisinage \deleted{tubulaire}} \inserted{\uline{conique}}
\inserted{(\uline{``absolu''})} de \deleted{$Y$}
\inserted{(ou ``cône sur $Y$'')} est un espace topologique
$\bar T\deleted{\unsure{_Y}}$ contenant $Y$
(la ``\uline{base}''), muni d'une partie fermée $\dot T$
(le ``bord'' du tube'') et d'une opération du monoïde topologique
$I=[0,1]$ (pour la multiplication) sur $\bar T$, notée
\[
(1)\qquad (\lambda,t)\longmapsto\lambda t
\qquad I\times\bar T\longrightarrow\bar T
\]
satisfaisant :
\[
\begin{array}{ll}
a)&\lambda y=y\quad\forall y\in Y\\
b)&0.\deleted{\unsure{\bar T}}\bar T\subset Y
\end{array}
\]
\margin{donc les $\lambda_{\bar T}:\bar T\longrightarrow\bar T$
$(0\leq\lambda\leq1)$ forment une homotopie de
l'identité $\operatorname{id}_{\bar T}$ (pour $\lambda=1$)
sur une \uline{rétraction} de $\bar T$ sur $Y$,
laissant les pts de $Y$ fixes au cours de l'homotopie}
\margin{\uline{NB} La condition b) \unsure{sous sa forme faible}
$b')\ 0.\dot T\subset Y$, plus la condition c)
(\unsure{à qui celle-ci donne un sens}) suffisent :
\unsure{entraînent} (b).}

c) Le \deleted{application} diagramme d'applications
\diagram{\[
(2)\qquad
\begin{tikzcd}[column sep=large]
I\times\dot T \arrow[r,"{(\lambda,\unsure{x})\mapsto\lambda x}",
"\deleted{\unsure{g}}"'] & \bar T\\
\{0\}\times\dot T \arrow[u,hook] \arrow[r,"{(0,x)\mapsto0.x}"'] & Y \arrow[u,hook,"\mathrm{inc}"']
\end{tikzcd}
\]}
est \uline{cocartésien}$^{*}$, i.e. identifie $\bar T$ au cône de l'application
\[
(3)\qquad f:\dot T\longrightarrow Y\qquad f(x)=0.x
\]
\margin{$^{*}$ Cela signifie que
$]0,1]\times\dot T\xrightarrow[\sim]{\mathrm{bij.}}\bar T\setminus Y$
et qu'une partie de $\bar T$ est ouverte ssi
\deleted{\unsure{son intersection}} \inserted{\unsure{ses images}}
\unsure{dans $Y$ et ses images dans $I\times\dot T$}.}
\[
(4)\qquad \bar T\simeq C(f)
\]
Dans cette identification, $Y\hookrightarrow\deleted{\unsure{T}}\bar T
\hookleftarrow\dot T$ correspondent
aux applications canoniques \deleted{\unsure{du cône $f$}}
$\operatorname{but}(f)\hookrightarrow C(f)\hookleftarrow
\operatorname{source}(f)$
et l'opération de $I$ sur $\bar T$ provient de l'opération de
$I$ sur $C(f)$, définie par passage au quotient de
l'opération évidente de $I$ sur $I\times\dot T$, et l'op. triviale
de $I$ sur $Y$. Les axiomes a) b) c) signifient
aussi que la situation $(\unsure{Y-\bar T\longrightarrow\dot T},
\unsure{I\times\bar T-\bar T})$
est isomorphe à celle d'un $C(f)$ avec $f:\dot T\longrightarrow Y$.
Notons que $\dot T$ est \unsure{connu quand on connaît} l'opération
de $I$ sur $\bar T$, \deleted{\unsure{comme on}} et posons
\[
(5)\qquad \mathring T=\bar T\setminus\dot T
\quad\text{(l'\uline{intérieur} du tube)}
\qquad\text{et notons que }\mathring T=\bigcup_{\lambda\in[0,1[}\lambda\bar T
\]
\margin{11}
